\subsubsection{Mã hóa File}\label{output-ma-hoa-file}

Trong hệ thống, sau khi các file JSON được sinh ra từ quá trình audit, việc bảo vệ tính toàn vẹn của các file này là rất quan trọng. Nếu file kết quả bị chỉnh sửa sau khi audit, báo cáo có thể không còn phản ánh đúng trạng thái thực tế của Docker Host. Vì vậy, hệ thống áp dụng cơ chế tạo mã băm SHA-256 để kiểm tra file có bị thay đổi hay không.

Cần phân biệt rõ rằng SHA-256 không phải là mã hóa theo nghĩa che giấu nội dung. Mã hóa là quá trình biến dữ liệu gốc thành dữ liệu không đọc được và có thể giải mã lại bằng khóa phù hợp. Trong khi đó, SHA-256 là hàm băm một chiều, dùng để tạo ra một chuỗi đại diện duy nhất cho nội dung dữ liệu. Từ mã băm không thể khôi phục lại dữ liệu gốc. Mục đích của SHA-256 là phát hiện sự thay đổi dữ liệu, không phải để giấu dữ liệu.

Trong đồ án, cơ chế hash được sử dụng để đảm bảo tính toàn vẹn của bằng chứng kiểm toán. Khi một file JSON được tạo ra, hệ thống có thể tính mã băm của toàn bộ file. Giá trị này được gọi là \texttt{source\_sha256}. Đây là dấu vân tay số của file JSON gốc.

Ví dụ lệnh kiểm tra hash của file JSON:

\begin{verbatim}
sha256sum audit/docker_benchmark_reports/192.168.1.8_section5.json
\end{verbatim}

Kết quả trả về là một chuỗi SHA-256. Nếu nội dung file JSON thay đổi dù chỉ một ký tự, giá trị SHA-256 sẽ thay đổi hoàn toàn. Nhờ đó, hệ thống có thể phát hiện file đã bị chỉnh sửa hoặc không còn nguyên vẹn so với thời điểm ban đầu.

Ngoài hash của file gốc, hệ thống còn có thể tính hash cho dữ liệu sau khi đã được chuẩn hóa và lưu vào Evidence Store. Giá trị này được gọi là \texttt{record\_sha256}. Nếu \texttt{source\_sha256} dùng để bảo vệ file JSON gốc, thì \texttt{record\_sha256} dùng để bảo vệ dữ liệu sau xử lý.

Quy trình tổng quát như sau:

\begin{enumerate}
\item Ansible chạy audit và sinh file JSON.
\item Hệ thống tính \texttt{source\_sha256} từ file JSON gốc.
\item Dữ liệu JSON được đọc và chuẩn hóa.
\item Hệ thống tính \texttt{record\_sha256} từ dữ liệu đã chuẩn hóa.
\item Các giá trị hash được lưu lại để phục vụ kiểm tra về sau.
\item Khi cần xác minh, hệ thống tính lại hash hiện tại và so sánh với hash đã lưu.
\end{enumerate}

Nếu hash hiện tại trùng với hash đã lưu, dữ liệu được xem là còn nguyên vẹn. Nếu hash không trùng, hệ thống có thể kết luận rằng dữ liệu đã bị thay đổi, bị lỗi hoặc không còn đáng tin cậy.

Cơ chế này đặc biệt quan trọng trong bài toán kiểm toán vì kết quả audit cần có độ tin cậy cao. Khi báo cáo được sử dụng làm bằng chứng, người kiểm tra cần đảm bảo dữ liệu trong báo cáo là dữ liệu thật được sinh ra từ quá trình audit, không phải dữ liệu đã bị chỉnh sửa thủ công.

Như vậy, phần ``Mã hóa File'' trong đồ án thực chất tập trung vào việc bảo vệ tính toàn vẹn của file thông qua mã băm SHA-256. Đây là cơ sở để hệ thống hỗ trợ nguyên tắc Integrity \& Hashing, giúp dữ liệu kiểm toán đáng tin cậy hơn và có thể xác minh lại khi cần.
